% ゲームプログラミングワークショップ 2026 投稿用 Extended Abstract
% 投稿規定: A4縦・2ページ以内・フォント10pt以上・PDF・EasyChair (GPW-26)
% ビルド: make -C paper abstract （lualatex を2回実行）
% 本稿は投稿用の extended abstract であり、本論文（main）とは別稿である。
\documentclass[10pt,a4paper,twocolumn]{ltjsarticle}
\usepackage[top=20mm,bottom=22mm,left=18mm,right=18mm]{geometry}
\usepackage{tikz}
\usepackage{booktabs}
\usepackage[hidelinks]{hyperref}
\usepackage{fancyhdr}
\setlength{\columnsep}{7mm}
\pagestyle{plain}

% 1ページ目の柱に投稿区分を出す（本文の版面は消費しない）
\setlength{\headheight}{14pt}
\fancypagestyle{gpwfirst}{%
  \fancyhf{}%
  \fancyhead[R]{\small ゲームプログラミングワークショップ 2026\quad Extended Abstract}%
  \fancyfoot[C]{\thepage}%
  \renewcommand{\headrulewidth}{0pt}%
}

% 盤面図（paper/tools/sfen2tikz.py が生成する noroshi_board.tex 用）
\newcommand{\koma}{\fontsize{13}{13}\selectfont}
\newcommand{\mochi}{\normalsize}

\begin{document}

\twocolumn[{%
\begin{center}
{\Large\bfseries 協力詰における全駒煙詰の発見}\\[1mm]
{\normalsize ---探索手法の逐次的な設計と改良による計算機支援創作---}\\[1.5mm]
{\small Discovery of a Full-Set (40-Piece) Smoke Mate in Helpmate Shogi:\\
Computer-Aided Composition through Iterative Design and Refinement of Search Methods}\\[3mm]
{\normalsize Keigo Oka\footnotemark[1]}
\end{center}
\vspace{1mm}
\begin{quote}\small
\textbf{概要}\quad
協力詰（ばか詰）において，将棋駒一組40枚すべてを盤上に置き持駒なしで始まり，
詰上がりが3枚となる「全駒煙詰」を計算機で探索した．
詰上がりからの逆算幅優先探索に，各候補局面の解一意性の完全検証，
探索記録から学習した評価関数によるビーム選択，および解集合の合流構造の解析に
基づく探索の再始動を組み合わせ，40枚・113手・解が一意の協力詰の煙詰作品を得た．
局面の完全性は完全探索により確認し，発表作は独立の検討ツールでも検証した．
40枚は初形駒数の理論上限である．これは網羅的探索ではなく、まったく異なる全駒煙詰が存在する可能性もある．
\end{quote}
\vspace{2mm}
}]
\thispagestyle{gpwfirst}
\footnotetext[1]{Google. Work done in a personal capacity（本研究は所属組織と無関係に個人として行った）．}
\setcounter{footnote}{1}

\section{はじめに}

協力詰（ばか詰）は，攻方と受方が協力して最短手数で受方玉を詰める
フェアリー詰将棋のルールであり，チェス・プロブレムの
helpmate（1854年，Lange~\cite{lange}）に対応する．
日本では1961年に『詰将棋パラダイス』67号でS・オギノ氏が「ばか詰」として提唱し，
1970年の中出慶一氏・加藤徹氏らの作品発表を経て，1971年の常設欄「ばか詰教室」
開設で定着した\cite{kanna}（「協力詰」は後年の言い換えである）．
解手順が一意であることが完全作の要件である．
煙詰は，初形で盤上に多数の駒を置き，詰上がりが玉を含む3枚のみとなる
詰将棋のジャンルで，伊藤看寿『将棋図巧』第99番（1755）\cite{zuko}に始まる．
\textbf{全駒煙}とは、初形が攻方玉を除く39枚，もしくは攻方玉を含む40枚すべての駒で始まる煙詰のことである。

協力詰では受方も詰みに協力するため玉は少数の駒で早く詰み，
多数の駒を一意な手順で消費し尽くすことが難しい．
2004年のウェブ上の議論\cite{tetsu2004}での到達は神無三郎氏の25枚までで，
「39枚（双玉なら40枚）から3枚にするのはおそらく不可能でしょう」と
評されていた\cite{mozu2004}．著者の知る限り，以後も全駒煙の作例は知られていない．
本稿では初形40枚・解が一意の協力詰の煙詰を計算機探索を交えた創作活動で
得たことを報告する．その1局（図\ref{fig:noroshi}）は『詰将棋パラダイス』
2026年8月号に特別懸賞出題される予定である\footnote{ogiekako「狼煙（のろし）」
協力詰113手．著者は筆名 ogiekako で作品発表を行っている．}．

\begin{figure}[t]
\centering
% generated by sfen2tikz.py -- do not edit
% sfen: 8+P/6K1p/4S3+P/3+Pp+PP+pn/2BpPP1p1/3Bk+p+p1P/1RRs+lgLNG/2SPNlNPG/G+pL4S1 b - 1
\begin{tikzpicture}[x=-4.8mm,y=-5.2mm]
\foreach \i in {1,...,8} {
  \draw[line width=0.1mm] (\i,0) -- (\i,9);
  \draw[line width=0.1mm] (0,\i) -- (9,\i);}
\draw[line width=0.4mm] (0,0) rectangle (9,9);
\foreach \p in {(3,3),(6,3),(3,6),(6,6)} \fill \p circle [radius=0.3mm];
\node[inner sep=0pt] at (0.5,0.5) {\koma と};
\node[inner sep=0pt] at (2.5,1.5) {\koma 王};
\node[inner sep=0pt,rotate=180] at (0.5,1.5) {\koma 歩};
\node[inner sep=0pt] at (4.5,2.5) {\koma 銀};
\node[inner sep=0pt] at (0.5,2.5) {\koma と};
\node[inner sep=0pt] at (5.5,3.5) {\koma と};
\node[inner sep=0pt,rotate=180] at (4.5,3.5) {\koma 歩};
\node[inner sep=0pt] at (3.5,3.5) {\koma と};
\node[inner sep=0pt] at (2.5,3.5) {\koma 歩};
\node[inner sep=0pt,rotate=180] at (1.5,3.5) {\koma と};
\node[inner sep=0pt,rotate=180] at (0.5,3.5) {\koma 桂};
\node[inner sep=0pt] at (6.5,4.5) {\koma 角};
\node[inner sep=0pt,rotate=180] at (5.5,4.5) {\koma 歩};
\node[inner sep=0pt] at (4.5,4.5) {\koma 歩};
\node[inner sep=0pt] at (3.5,4.5) {\koma 歩};
\node[inner sep=0pt,rotate=180] at (1.5,4.5) {\koma 歩};
\node[inner sep=0pt] at (5.5,5.5) {\koma 角};
\node[inner sep=0pt,rotate=180] at (4.5,5.5) {\koma 玉};
\node[inner sep=0pt,rotate=180] at (3.5,5.5) {\koma と};
\node[inner sep=0pt,rotate=180] at (2.5,5.5) {\koma と};
\node[inner sep=0pt] at (0.5,5.5) {\koma 歩};
\node[inner sep=0pt] at (7.5,6.5) {\koma 飛};
\node[inner sep=0pt] at (6.5,6.5) {\koma 飛};
\node[inner sep=0pt,rotate=180] at (5.5,6.5) {\koma 銀};
\node[inner sep=0pt,rotate=180] at (4.5,6.5) {\koma 杏};
\node[inner sep=0pt,rotate=180] at (3.5,6.5) {\koma 金};
\node[inner sep=0pt] at (2.5,6.5) {\koma 香};
\node[inner sep=0pt] at (1.5,6.5) {\koma 桂};
\node[inner sep=0pt] at (0.5,6.5) {\koma 金};
\node[inner sep=0pt] at (6.5,7.5) {\koma 銀};
\node[inner sep=0pt] at (5.5,7.5) {\koma 歩};
\node[inner sep=0pt] at (4.5,7.5) {\koma 桂};
\node[inner sep=0pt,rotate=180] at (3.5,7.5) {\koma 香};
\node[inner sep=0pt] at (2.5,7.5) {\koma 桂};
\node[inner sep=0pt] at (1.5,7.5) {\koma 歩};
\node[inner sep=0pt] at (0.5,7.5) {\koma 金};
\node[inner sep=0pt] at (8.5,8.5) {\koma 金};
\node[inner sep=0pt,rotate=180] at (7.5,8.5) {\koma と};
\node[inner sep=0pt] at (6.5,8.5) {\koma 香};
\node[inner sep=0pt] at (1.5,8.5) {\koma 銀};
\node[inner sep=0pt] at (9.75,0.4) {\mochi ☖};
\node[inner sep=0pt] at (9.75,1.22) {\mochi な};
\node[inner sep=0pt] at (9.75,2.04) {\mochi し};
\node[inner sep=0pt] at (-0.75,0.4) {\mochi ☗};
\node[inner sep=0pt] at (-0.75,1.22) {\mochi な};
\node[inner sep=0pt] at (-0.75,2.04) {\mochi し};
\end{tikzpicture}

\caption{「狼煙」協力詰113手（盤上40枚・詰上がり3枚）．}
\label{fig:noroshi}
\end{figure}

\section{逆算探索と一意性検証}

基盤には著者が開発したRust製の協力詰ソルバー fmrs~\cite{fmrs} を用いた．
fmrs は既知作品による回帰テストを備える
（加藤徹氏作19447手の「寿限無」（1975）を約5秒で解く）．
協力詰には攻防の対立構造がないため，普通詰将棋の解図で
主力となる証明数探索（df-pn等）は前提を欠く．完全性（余詰なし）の判定は
「最短で詰みに至る手順がちょうど1本」という解経路数の条件として扱う．

探索は，詰将棋創作で用いられる逆算法\cite{hirose1998,hirose2026}の計算機化である．
3枚の詰上がりを起点に，層別の幅優先探索で1手ずつ逆向きに展開し，
生成された各候補局面について解の一意性をキャッシュ付きの深さ制限DFSで完全検証する．
逆算のフロンティアは2手あたり約1.7--1.9倍で増加し，40枚に必要な
70手超の遡行では$10^{10}$局面を超える規模になる．実測では768\,GBメモリの
計算機による全幅探索で22枚が上限だった（表\ref{tab:results}）．
全幅探索の高速化だけで必要規模に到達する見込みは薄いため，
以降は候補選択法（ビーム化）を検討した．
なお本探索では歩の枚数比率や飛角・香桂の許容枚数などの制約を課しており，
以降の結果はこの制約クラスの下でのものである．<- TODO: ここ、実際どういう制限で例えばやったのか、commit msg とかも参考にして記述したほうがいい。

\section{学習誘導ビーム探索}

学習モデルによる探索誘導（cf.~\cite{alphago}）を採る．
全幅探索が可能な範囲を解析すると，より深い最大駒数解に寄与する局面は
フロンティアのごく一部であり，その割合は深いほど下がる
（詰みから11手で2.6\%，31手で0.02\%）．これを取り落とさない選択規則が
必要になる．そこで局面の特徴量85個（玉の可動域，駒の分布，成駒数など）から
「その局面から到達可能な最大駒数」を回帰する勾配ブースティング木を学習した．
ラベルは，探索で生成した局面の親子関係の全記録から，後ろ向きの動的計画法で
厳密に計算した．評価は解系列単位でデータを分割して行い（同一系列の学習・検証への混入を防ぐ），
同一（深さ，駒数）内の順位相関はラベルの厳密化で0.26から0.49に上がった．

選択規則の寄与も大きく，スコア上位の貪欲選択は類似局面に偏って悪化するため，
温度付き乱択（Gumbel top-$K$）と駒数別の枠確保を併用した．
ビーム幅$10^5$での到達駒数は乱択の30枚に対し誘導ビームで35枚，
幅$3\times10^5$で37枚，さらに幅を深さに応じて増やす設定で39枚となった．
ただし到達駒数の実行間分散は$\pm$3--6枚と大きく，
表\ref{tab:results}は少数回の実行による到達値であって統制された比較ではない．

\begin{table}[t]
\centering
\caption{各探索設定の到達値（全幅探索の行のみ制約クラスが異なる）．}
\label{tab:results}
\small
\begin{tabular}{lrrl}
\toprule
探索設定 & 駒数 & 手数 & 備考 \\
\midrule
全幅探索 & 22 & --- & 768\,GBで頭打ち \\
乱択ビーム & 30 & --- & 幅$10^5$ \\
学習誘導ビーム & 35 & --- & 同幅（$3{\times}10^5$で37） \\
〃＋幅ランプ & 39 & 107 & 438局 \\
合流点からの再探索 & 40 & 113 & 2618局 \\
\bottomrule
\end{tabular}
\end{table}

TODO: 幅ランプ、って耳慣れない単語じゃない？

\section{合流構造の解析と40枚}

手動リスタートをいくつかの異なる幅ランプで行った結果，初形39枚の解が得られた。のこり1枚の歩を詰め込めるといいところまできた。
で、39枚で頭打ちになった探索の解集合を解析した．TODO: ここ、口語的すぎる

39枚438局の解を詰み側から辿ると，詰みから82手までは全解が同一の手順を通り，
分岐は初形側の約25手に限られていた．全解が通る最深の共通局面（以下\textbf{合流点}）
を逆算の起点に据えると，それ以深は局面数が少なく全幅探索で数十秒で調べ尽くせる．この方法で，39枚の互いに独立な3系列のそれぞれについて，
「その合流点以深に40枚は存在しない」ことを確認した．
ただし，この不存在が示されるのは各合流点以深の範囲に限られる．
したがって40枚が在るとすれば，より詰みに近い位置で分岐する系列に限られる．

そこで3系列すべての最深共通局面（詰みから42手・14枚）を起点とし，
ビーム幅を深部ほど増やして（最大$10^7$）再探索した（96 vCPU・768\,GB）．
TODO: これは、GCP の spot instance を短期的に借りてやった。
このビーム探索は候補フロンティアが空になるまで継続し，
詰みから113手の深さで40枚・持駒なしの局面2618局が得られた．

\section{検証と作品化}

得られた2618局のそれぞれについて，協力手順が一意であることを
fmrs の完全探索により確認した．ビームは候補生成の絞り込みにのみ用い，
完全性判定には関与しない．発表作については，
神無次郎氏による独立のフェアリー詰将棋検討・創作支援ツール fm~\cite{fm}
でも完全性を確認した．一方で，今回発見した局面で全駒煙詰のすべてであることは保証されない
（枝刈りされた範囲に別の全駒煙詰が存在するであろう）．

作品化では，成駒の種類・枚数や攻方駒の玉への密集度などで2618局を
順位付けして3局に絞り，さらにそれを、自陣成り駒の少なさといった、人間の美的感覚に基づき微調整を行った
（修正後も解の一意性を再検証）．

\section{おわりに}

初形駒数の上で40枚は上限である．一方で、今回探索した領域は可能な問題空間のごく一部に過ぎず、本作品は協力詰の煙詰の可能性の端緒を示すものに過ぎない。
学習した評価関数による逆算探索の誘導は，他のフェアリールールや
チェス・プロブレムでの記録探索にも適用しうる．普通詰将棋では煙詰の自動生成が
報告されており\cite{hirose2026}，本報告はルールの異なる協力詰側から
これを補完する事例である．

TODO: 謝辞として、神無七郎氏：詰将棋パラダイスにおける特別出題の担当者。また、私が進捗報告を送るたびに励ましをいただいた。

\begin{thebibliography}{9}
\small
\bibitem{zuko} 伊藤看寿：『将棋図巧』第99番，1755．
\bibitem{lange} M. Lange: \textit{Deutsche Schachzeitung}, Dec.\ 1854.
\bibitem{kanna} 神無七郎：神無一族の氾濫 第19回．
\url{k7ro.sakura.ne.jp/overflow/hr19.htm}
\bibitem{tetsu2004} TETSU：協力詰の煙詰．詰将棋メモ，2004．
\url{toybox.tea-nifty.com/memo/2004/10/post_23.html}
\bibitem{mozu2004} mozuyama：ばか煙．勝手に将棋トピックス，2004．
\url{mozuyama.hatenadiary.org/entry/20041024/P20041024BAKAKEMURI}
\bibitem{hirose1998} 広瀬正幸，伊藤琢巳，松原仁：逆算法による詰め将棋の自動創作．
人工知能学会誌，Vol.~13，No.~3，pp.~452--460，1998．
\bibitem{hirose2026} 広瀬稔，馬屋原剛：独自の評価関数に基づいた煙詰の自動生成．
情報処理学会研究報告，Vol.~2026-GI-57，No.~4，2026．
\bibitem{alphago} D. Silver, et al.: Mastering the game of Go with deep neural
networks and tree search. \textit{Nature}, Vol.~529, pp.~484--489, 2016.
\bibitem{fmrs} K. Oka: fmrs．\url{github.com/ogiekako/fmrs}
\bibitem{fm} 神無次郎：fm．\url{www.dokidoki.ne.jp/home2/takuji/FM.html}
\end{thebibliography}

\end{document}
